The broad field summaries mix poorly: their rank-normalized $\widehat R$ values range from 1.63 to 2.25, with bulk ESS estimates between 5.0 and 6.7 (Table~\ref{tab:field-mixing}). For the survey-footprint mean, the average within-chain autocorrelation is 0.999 at lag one and 0.896 at lag 100. The traces in Figure~\ref{fig:field-mixing} show persistent drift and disagreement, so these small ESS values are diagnostic warnings rather than a basis for precise Monte Carlo uncertainty estimates.

Local coordinates also show substantial dependence: among the 128 monitored footprint vertices, the median $\widehat R$ is 1.136, the maximum is 1.564, and 127 exceed 1.01; median bulk ESS is 21.7. Outside the footprint, 31 of the 32 monitored vertices exceed 1.01. Consequently, the field diagnostics support the earlier warning about incomplete posterior exploration, including spatial averages that can mix much more slowly than an individual coordinate. All fields were initialized at zero, so this check does not establish robustness to dispersed field starts.

The small conditional blocks provide a plausible mechanism for slow movement on large spatial scales, while dependence between the field, site counts, breeding coefficients and dispersal width may also contribute. These runs do not separate those causes or compare alternative samplers; larger-scale or joint updates would require their own empirical mixing assessment.
